\documentclass[12pt]{article}
\usepackage[margin=1in]{geometry}
\usepackage{enumitem}
\usepackage{titlesec}
\usepackage{parskip}
\usepackage{graphicx}
\usepackage{hyperref}
\usepackage{fontenc}

\title{\textbf{Selznick 1949 RG Notes}\\
\large Selznick, Philip. \textit{TVA and the Grass Roots: A Study in the Sociology of Formal Organization}.\\
Berkeley: University of California Press, 1949.}
\author{}
\date{}

\begin{document}

\maketitle

\section{Main Idea}



%--------------------------------------------------
\section{General Notes}
%--------------------------------------------------

\begin{itemize}[leftmargin=*, itemsep=4pt]
    \item Organizations are faced with an existential gap between what they wish to do and what they are able to accomplish; subsequently, the natural outome is ``a reconciliation between the desire and the ability.''
    \item Organizations ``take on the aspects of a social organism,'' with the following consequences:
    \begin{enumerate}
        \item Decentralization without harming organizational goals is possible only when a common set of beliefs are deeply held / indoctrinated amongst multiple levels of leadership \emph{below the very top} -- that is, organizational goals and values are internalized such that a greater degree of discretion is possible
        \item Consistent doctrine allows for technologically-specialized experts to bind together when they may otherwise be quite drawn to their individual groups / professional interests. 
        \item Institutions gain a tendency to defend its own functions and existence. ``These defensive activities are aided when a set of beliefs is so fashioned as at once to fortify the special needs or interests of the organization and to provide an aura of disinterestedness under which formal discussion may be pursued.''
    \end{enumerate}
    \item Institutions must adapt to their environments, but: at least in the case of the TVA, adaptation to local institutions was much more significant (both in presence and necessity) than adaptation to local populations. Though local populations may very well hold strong beliefs, it is the institutions which represent those populations who may actually impact the route to implementation, positively or negatively. (For the TVA, this apporoach represented a pragmatic rather than normative effort -- grassroots organizations allowed for effective progress, not simply a reflection of ``good governance.'')
    \item Eight inherent dilemmas regarding administrative/policymaking discretion (all quotes):
    \begin{enumerate}
        \item There is a dilemma of doctrine and commitment, or of the abstract and concrete.
        \item There is also the dilemma of consent and conformance, or of selective decision and total involvement.
        \item So far as discretion is delegated, bifurcation of policy and administration is reinforced.
        \item Theories about government become preempted by social forces.
        \item Emphasis on existing institutions as democratic instruments may wed the agency to the status quo.
        \item Decision at the grass roots may be inhibited by the system of national pressure groups.
        \item Commitment to existing agencies may shape and inhibit policy in unanticipated ways.
        \item Existing agencies inhibit a direct approach to the local citizenry.
    \end{enumerate} 
\end{itemize}


\section{Important Specific Factual Claims}

\begin{itemize}[leftmargin=*, itemsep=4pt]
    \item ``Organizations established in a stable context, with assured futures, do not feel the same urgencies as those born of turmoil and set down to fend for themselves in undefined ways among institutions fearful and resistant. As in individuals, insecurity summons ideological reinforcements.'' Broader points of institutional threat -- argument possibly that it could result in better institutions under such stressors, though not entirely clear
\end{itemize}


\section{Connections}

\begin{itemize}
    \item Discussions of cooptation and decentralization are largely aligned with Huber \& Shipan -- ideological congruency allows for greater discretion, which in turn is productive for facilitating grassroots organization. A broader tension though arises when there is not ideological congruency and therefore cannot be sufficient discretion given to lower levels of authority in order to promote a grassroots approach.
    \item Institutions adapting to existing institutions rather than existing populations -- very much a Lijphart/Powell viewpoint of representative democracy (majoritarian) whereby elected officials are granted a mandate
\end{itemize}


\section{My Critiques}

\begin{itemize}
    \item
\end{itemize}


\end{document}
